\documentclass[11pt,a4paper]{article}

% ===================== PAQUETES =====================
\usepackage[utf8]{inputenc}
\usepackage[spanish]{babel}
\usepackage[T1]{fontenc}
\usepackage[margin=2.5cm]{geometry}
\usepackage{graphicx}
\usepackage{tabularx}
\usepackage{array}
\usepackage{xcolor}
\usepackage{colortbl}
\usepackage{float}
\usepackage{verbatim}
\usepackage{hyperref}

\renewcommand{\arraystretch}{1.3}

% =========================================================
% SISTEMA DE CARPETAS PARA LAS IMÁGENES
% =========================================================
% Todas las imágenes (PNG) van dentro de la carpeta "figuras/":
%
%   figuras/logo/LogoUni.png
%   figuras/casos_uso/diagrama_casos_uso.png
%   figuras/secuencia/CU-C01_secuencia.png ... CU-C07_secuencia.png
%   figuras/comunicacion/CU-C01_comunicacion.png ... CU-C07_comunicacion.png
%
% Basta con subir cada PNG a su carpeta con el nombre indicado; no hay
% que modificar este archivo. Si una imagen aún no existe, el documento
% compila igual y muestra un recuadro con la ruta esperada.
\graphicspath{{./}{figuras/}}

% ===================== LOGO (opcional) =====================
\newcommand{\LogoUni}{%
    \IfFileExists{figuras/logo/LogoUni.png}{%
        \includegraphics[width=0.35\textwidth]{figuras/logo/LogoUni.png}\\[1cm]%
    }{}%
}

% ===================== COMANDO PARA TABLA DE CASO DE USO =====================
% Uso:
% \CasoUso{Nombre}{ID}{Actores}{Descripcion}{Precondicion}{Primario}{Alternativo}{Excepcional}

\newcommand{\CasoUso}[8]{
\begin{center}
\begin{tabularx}{\textwidth}{|l|X|}
\hline
\rowcolor{gray!25}
\multicolumn{2}{|c|}{\textbf{Caso de Uso: #1}} \\
\hline
\textbf{ID} & #2 \\
\hline
\textbf{Actor(es)} & #3 \\
\hline
\textbf{Descripción} & #4 \\
\hline
\textbf{Precondición} & #5 \\
\hline
\rowcolor{gray!15}
\multicolumn{2}{|c|}{\textbf{Escenarios}} \\
\hline
\textbf{Principal} & #6 \\
\hline
\textbf{Alternativo(s)} & #7 \\
\hline
\textbf{Excepcional(es)} & #8 \\
\hline
\end{tabularx}
\end{center}
\vspace{0.5cm}
}

% ===================== COMANDO PARA INSERTAR DIAGRAMAS =====================
% Uso: \Diagrama{ruta/imagen.png}{etiqueta}{Leyenda}
% Si la imagen todavía no existe, se muestra un recuadro de espacio reservado.
\newcommand{\Diagrama}[3]{
\begin{figure}[H]
    \centering
    \IfFileExists{#1}{%
        \includegraphics[width=\textwidth,height=0.72\textheight,keepaspectratio]{#1}%
    }{%
        \fbox{\parbox[c][0.35\textheight][c]{0.85\textwidth}{\centering
        \textit{Insertar aquí la imagen del diagrama}\\[0.4em]
        \texttt{\detokenize{#1}}}}%
    }
    \caption{#3}
    \label{#2}
\end{figure}
}

% ===================== COMANDO PARA FICHA DE DIAGRAMA =====================
% Uso: \Ficha{ID}{Caso de uso}{Participantes}{Escenarios representados}{Elementos de UML usados}
\newcommand{\Ficha}[5]{
\begin{center}
\begin{tabularx}{\textwidth}{|l|X|}
\hline
\rowcolor{gray!25}
\multicolumn{2}{|c|}{\textbf{#1 -- #2}} \\
\hline
\textbf{Participantes} & #3 \\
\hline
\textbf{Escenarios representados} & #4 \\
\hline
\textbf{Elementos UML} & #5 \\
\hline
\end{tabularx}
\end{center}
}

% ===================== PORTADA =====================

\title{

    \LogoUni
    \textbf{Análisis y Especificación de Casos de Uso\\
    y Diagramas de Interacción}\\
    \vspace{0.5cm}
    \large Tinder UD
    \vspace{0.3cm}
}

\author{
    \vspace{0.7cm}
    \textbf{Autores:}\\
    Angel Andres Diaz Vergara 20231020056\\ Gabriel Fernando Lozano Echeverry - 20232020105\\
    Juan Camilo Mosquera Palomino - 20241020120\\[1.5em]
    \vspace{0.7cm}
    \textbf{Profesor:}\\
    Sandro Javier Bolaños Castro\\[1.5em]
    \vspace{0.7cm}
    \textbf{Asignatura:}\\
    Fundamentos de ingeniería de software\\[1.5em]
    \vspace{0.7cm}
    \textbf{Universidad:}\\
    Universidad Distrital Francisco José de Caldas\\
    Facultad de Ingeniería\\
    Bogotá D.C., Colombia\\[1em]
}
\date{\today}

\begin{document}

\maketitle

\tableofcontents

\newpage

% =========================================================
% PARTE I. ANÁLISIS Y ESPECIFICACIÓN DE CASOS DE USO
% =========================================================

\part{Análisis y Especificación de Casos de Uso}

% ---------------------------------------------------------
% 1. PLANTEAMIENTO
% ---------------------------------------------------------

\section{Planteamiento del Software}

La propuesta consiste en desarrollar una aplicación de interacción social orientada
exclusivamente a estudiantes de una comunidad universitaria. La aplicación busca
facilitar que los estudiantes conozcan otras personas pertenecientes a la misma
universidad y encuentren posibles afinidades a partir de sus intereses y de la
información que cada usuario decida compartir.

El sistema permitirá que cada estudiante cree un perfil personal en el que pueda
agregar fotografías, una biografía, etiquetas relacionadas con sus intereses y
información académica básica, como la facultad y la sede a la que pertenece.

A partir de estos perfiles, los estudiantes podrán explorar otras personas y
encontrar aquellas que compartan intereses o características de su preferencia.
Cuando exista interés mutuo entre dos estudiantes, se establecerá una conexión
(\textit{match}), después de la cual se habilitará un canal de comunicación mediante
un chat.

La aplicación también contempla mecanismos básicos de seguridad para que los
usuarios puedan bloquear o reportar a otros estudiantes cuando consideren que una
interacción no es apropiada.

% ---------------------------------------------------------
% 2. PROPÓSITO
% ---------------------------------------------------------

\section{Propósito}

El propósito principal de la aplicación es facilitar la conexión entre estudiantes
de una misma comunidad universitaria, proporcionando un espacio en el que puedan
conocerse, identificar intereses en común y establecer comunicación cuando exista
interés mutuo.

El sistema busca que el proceso de conocer nuevas personas sea sencillo y esté
centrado en la información que los propios estudiantes decidan compartir sobre sí
mismos.

Para ello, la aplicación permitirá crear perfiles personales, explorar perfiles
de otros estudiantes, encontrar personas con intereses afines, establecer
conexiones y comunicarse mediante un chat.

Además, se busca proporcionar mecanismos básicos que permitan al usuario controlar
sus interacciones mediante las opciones de bloqueo y reporte.

% ---------------------------------------------------------
% 3. JUSTIFICACIÓN
% ---------------------------------------------------------

\section{Justificación}

Dentro de una comunidad universitaria existe una gran cantidad de estudiantes que
pertenecen a diferentes facultades, programas académicos y sedes. Aunque todos
forman parte de una misma comunidad, las actividades académicas y la distribución
de los estudiantes pueden dificultar que algunas personas conozcan a otras con
intereses similares.

La aplicación propuesta busca brindar una alternativa para facilitar este proceso,
permitiendo que los estudiantes conozcan personas de la misma universidad a partir
de información básica como sus intereses, biografía, fotografías, facultad y sede.

A diferencia de una aplicación de citas general, el sistema estará orientado a una
comunidad universitaria específica. Por esta razón, no se considera necesario
utilizar la ubicación geográfica como criterio principal. En su lugar, la
pertenencia a la universidad y la información relacionada con la facultad o sede
permitirán establecer el contexto de los usuarios.

La propuesta también contempla funciones de seguridad básicas, como bloquear y
reportar usuarios, con el objetivo de proporcionar mecanismos para controlar
interacciones no deseadas.

% ---------------------------------------------------------
% 4. ALCANCE
% ---------------------------------------------------------

\section{Alcance del Sistema}

El sistema estará orientado a estudiantes pertenecientes a la comunidad
universitaria.

Dentro del alcance inicial de la aplicación se encuentran las siguientes
funcionalidades:

\begin{itemize}
    \item Creación y gestión de una cuenta de usuario.
    \item Creación y edición de un perfil personal.
    \item Carga y eliminación de fotografías.
    \item Registro de una biografía.
    \item Registro de intereses mediante etiquetas.
    \item Información relacionada con la facultad y sede del estudiante.
    \item Exploración de perfiles de otros estudiantes.
    \item Búsqueda de personas con intereses afines.
    \item Aplicación de filtros relacionados con los perfiles.
    \item Expresión de interés hacia otro estudiante.
    \item Establecimiento de una conexión cuando existe interés mutuo.
    \item Comunicación mediante un chat después de establecer una conexión.
    \item Bloqueo de otros usuarios.
    \item Reporte de otros usuarios.
\end{itemize}

La propuesta no contempla inicialmente funcionalidades relacionadas con la
ubicación geográfica precisa de los usuarios, sistemas de pago, publicidad,
historias, contenido multimedia avanzado ni mecanismos de recomendación basados
en algoritmos complejos.

% ---------------------------------------------------------
% 5. ACTOR PRINCIPAL
% ---------------------------------------------------------

\section{Actor Principal}

\subsection{Estudiante}

El estudiante es el actor principal del sistema. Es la persona que utiliza la
aplicación para conocer otros estudiantes, compartir información sobre sí mismo,
buscar personas afines, establecer conexiones y comunicarse con aquellas personas
con las que haya establecido un \textit{match}.

El estudiante también puede administrar su cuenta y utilizar las funciones de
seguridad disponibles para bloquear o reportar a otros usuarios.

% ---------------------------------------------------------
% 6. IDENTIFICACIÓN DE CASOS DE USO
% ---------------------------------------------------------

\section{Identificación de Casos de Uso}

Los casos de uso se identifican desde una perspectiva orientada a los objetivos
del usuario y a la finalidad general del sistema.

Por lo tanto, los casos de uso no representan operaciones técnicas específicas de
la aplicación, sino las necesidades y objetivos que el estudiante busca alcanzar
mediante el uso del sistema.

El caso de uso principal identificado para la aplicación es
\textit{Facilitar la conexión entre estudiantes}. Los demás casos de uso
representan objetivos que permiten alcanzar o complementar dicha finalidad.

% ---------------------------------------------------------
% 7. TABLA GENERAL DE CASOS DE USO
% ---------------------------------------------------------

\section{Casos de Uso Identificados}

\begin{center}
\begin{tabularx}{\textwidth}{|c|X|X|X|}
\hline
\rowcolor{gray!25}
\textbf{ID} &
\textbf{Nombre del Caso de Uso} &
\textbf{Actor Principal} &
\textbf{Descripción} \\
\hline

CU-C01 &
\textbf{Facilitar la conexión entre estudiantes} &
Estudiante &
Permitir que los estudiantes conozcan personas de su comunidad universitaria y establezcan conexiones con aquellas que consideren afines. \\
\hline

CU-C02 &
\textbf{Presentarse ante otros estudiantes} &
Estudiante &
Permitir que el estudiante comparta información sobre sí mismo, como fotografías, biografía, intereses, facultad y sede. \\
\hline

CU-C03 &
\textbf{Encontrar personas afines} &
Estudiante &
Permitir que el estudiante explore y encuentre otros estudiantes de acuerdo con sus intereses y características de preferencia. \\
\hline

CU-C04 &
\textbf{Establecer una conexión} &
Estudiante &
Permitir que dos estudiantes establezcan una conexión cuando existe interés mutuo. \\
\hline

CU-C05 &
\textbf{Comunicarse con una conexión} &
Estudiante &
Permitir que los estudiantes que han establecido una conexión se comuniquen mediante un chat. \\
\hline

CU-C06 &
\textbf{Proteger la interacción} &
Estudiante &
Permitir que el estudiante bloquee o reporte a otros usuarios para controlar interacciones no deseadas o inapropiadas. \\
\hline

CU-C07 &
\textbf{Gestionar el acceso a la aplicación} &
Estudiante &
Permitir que el estudiante cree, acceda y gestione su cuenta dentro de la aplicación. \\
\hline

\end{tabularx}
\end{center}

% ---------------------------------------------------------
% 8. DIAGRAMA DE CASOS DE USO
% ---------------------------------------------------------

\section{Diagrama de Casos de Uso}

El diagrama de casos de uso representa la relación entre el actor principal y los
objetivos identificados para el sistema.

\Diagrama{figuras/casos_uso/diagrama_casos_uso.png}{fig:casos_uso}{Diagrama de casos de uso de Tinder UD}

% ---------------------------------------------------------
% 9. DESCRIPCIÓN DETALLADA DE CASOS DE USO
% ---------------------------------------------------------

\section{Descripción de Casos de Uso}

A continuación se describen los casos de uso desde la perspectiva del estudiante,
manteniendo el enfoque en los objetivos que busca alcanzar mediante la aplicación.

% ---------- CU-C01 ----------

\CasoUso
{Facilitar la conexión entre estudiantes}
{CU-C01}
{Estudiante}
{El estudiante utiliza la aplicación con el propósito de conocer otros estudiantes
de la comunidad universitaria y establecer conexiones con personas que considere
afines.}
{El estudiante debe pertenecer a la comunidad universitaria y tener acceso a la
aplicación.}
{
1. El estudiante accede a la aplicación. \newline
2. Consulta información de otros estudiantes. \newline
3. Encuentra personas que considera afines a sus intereses. \newline
4. Expresa interés en establecer una conexión. \newline
5. Cuando existe interés mutuo, se establece una conexión. \newline
6. El estudiante puede comunicarse con la persona conectada.
}
{
1. El estudiante puede continuar explorando perfiles sin establecer una conexión.
\newline
2. Puede modificar sus preferencias para encontrar otros perfiles.
}
{
1. No existen perfiles que coincidan con las preferencias del estudiante. \newline
2. El perfil de otro estudiante deja de estar disponible. \newline
3. El estudiante no puede completar la interacción.
}

% ---------- CU-C02 ----------

\CasoUso
{Presentarse ante otros estudiantes}
{CU-C02}
{Estudiante}
{El estudiante comparte información sobre sí mismo para permitir que otros usuarios
conozcan sus características, intereses y aspectos que considere relevantes.}
{El estudiante debe tener una cuenta dentro de la aplicación.}
{
1. El estudiante accede a su perfil. \newline
2. Agrega una o varias fotografías. \newline
3. Escribe una biografía. \newline
4. Selecciona etiquetas relacionadas con sus intereses. \newline
5. Indica información relacionada con su facultad y sede. \newline
6. Guarda la información del perfil.
}
{
1. El estudiante puede modificar posteriormente la información de su perfil. \newline
2. Puede eliminar fotografías o modificar sus intereses. \newline
3. Puede decidir qué información desea mantener en su perfil.
}
{
1. El estudiante intenta guardar información incompleta cuando existen datos
requeridos. \newline
2. Una fotografía no puede ser agregada al perfil. \newline
3. La información proporcionada no es válida.
}

% ---------- CU-C03 ----------

\CasoUso
{Encontrar personas afines}
{CU-C03}
{Estudiante}
{El estudiante busca otros estudiantes que puedan ser de su interés a partir de la
información disponible en sus perfiles y de sus preferencias.}
{El estudiante debe haber accedido a la aplicación y contar con un perfil.}
{
1. El estudiante accede a la sección de exploración. \newline
2. Consulta los perfiles disponibles. \newline
3. Utiliza sus preferencias para encontrar personas de interés. \newline
4. Puede utilizar criterios como intereses, facultad o sede. \newline
5. Consulta los perfiles que cumplen con sus criterios. \newline
6. Decide si desea establecer una conexión con alguno de ellos.
}
{
1. El estudiante puede realizar la búsqueda sin aplicar filtros específicos. \newline
2. Puede modificar sus criterios para obtener otros resultados. \newline
3. Puede continuar explorando perfiles diferentes.
}
{
1. No existen perfiles que coincidan con los criterios seleccionados. \newline
2. Los perfiles disponibles no contienen información suficiente para aplicar
determinados criterios.
}

% ---------- CU-C04 ----------

\CasoUso
{Establecer una conexión}
{CU-C04}
{Estudiante}
{El estudiante establece una conexión con otro estudiante cuando ambos expresan
interés de manera mutua.}
{El estudiante debe haber encontrado un perfil con el que desea establecer una
posible conexión.}
{
1. El estudiante expresa interés en otro perfil. \newline
2. El otro estudiante también expresa interés. \newline
3. Se identifica que existe interés mutuo. \newline
4. Se establece una conexión entre ambos estudiantes. \newline
5. Se habilita la posibilidad de comunicación entre los estudiantes conectados.
}
{
1. El estudiante puede continuar explorando otros perfiles después de expresar
interés. \newline
2. El estudiante puede establecer conexiones con diferentes usuarios.
}
{
1. No existe interés mutuo. \newline
2. Uno de los perfiles deja de estar disponible antes de establecer la conexión.
}

% ---------- CU-C05 ----------

\CasoUso
{Comunicarse con una conexión}
{CU-C05}
{Estudiante}
{El estudiante se comunica mediante un chat con otro estudiante con quien ha
establecido una conexión.}
{Debe existir una conexión entre ambos estudiantes.}
{
1. El estudiante consulta sus conexiones. \newline
2. Selecciona una conexión. \newline
3. Accede al espacio de conversación. \newline
4. Escribe un mensaje. \newline
5. Envía el mensaje. \newline
6. Puede recibir y consultar respuestas de la otra persona.
}
{
1. El estudiante puede consultar conversaciones anteriores. \newline
2. Puede abandonar la conversación y regresar posteriormente. \newline
3. Puede continuar la conversación en diferentes momentos.
}
{
1. La conexión deja de estar disponible. \newline
2. Uno de los usuarios bloquea al otro. \newline
3. El mensaje no puede enviarse temporalmente.
}

% ---------- CU-C06 ----------

\CasoUso
{Proteger la interacción}
{CU-C06}
{Estudiante}
{El estudiante controla sus interacciones con otros usuarios mediante las opciones
de bloqueo y reporte disponibles en la aplicación.}
{El estudiante debe haber accedido a la aplicación y debe poder identificar al
usuario con quien desea controlar la interacción.}
{
1. El estudiante accede al perfil o conversación de otro usuario. \newline
2. Selecciona la opción correspondiente para bloquear o reportar. \newline
3. Si realiza un reporte, indica el motivo. \newline
4. Confirma la acción. \newline
5. El sistema aplica la acción correspondiente.
}
{
1. El estudiante puede bloquear a un usuario sin realizar un reporte. \newline
2. Puede realizar un reporte y posteriormente bloquear al usuario. \newline
3. Puede cancelar la acción antes de confirmarla.
}
{
1. El usuario seleccionado deja de estar disponible. \newline
2. El reporte no contiene la información requerida. \newline
3. La acción no puede completarse temporalmente.
}

% ---------- CU-C07 ----------

\CasoUso
{Gestionar el acceso a la aplicación}
{CU-C07}
{Estudiante}
{El estudiante administra su acceso a la aplicación mediante la creación, ingreso
y gestión de su cuenta.}
{El estudiante debe pertenecer a la comunidad universitaria para utilizar la
aplicación.}
{
1. El estudiante selecciona la opción para registrarse. \newline
2. Proporciona la información requerida. \newline
3. Crea su cuenta. \newline
4. Puede iniciar sesión posteriormente. \newline
5. Puede utilizar la aplicación mientras su cuenta se encuentre activa.
}
{
1. El estudiante puede cerrar sesión y volver a acceder posteriormente. \newline
2. Puede modificar información relacionada con su cuenta. \newline
3. Puede solicitar la eliminación de su cuenta.
}
{
1. La información proporcionada durante el registro no es válida. \newline
2. La cuenta ya existe. \newline
3. El acceso no puede realizarse con la información proporcionada.
}

\newpage

% =========================================================
% PARTE II. DIAGRAMAS DE SECUENCIA Y DE COMUNICACIÓN
% =========================================================

\part{Diagramas de Secuencia y de Comunicación}

% ---------------------------------------------------------
% 10. INTRODUCCIÓN A LOS DIAGRAMAS
% ---------------------------------------------------------

\section{Introducción a los Diagramas de Interacción}

Esta parte del documento presenta los diagramas de secuencia y de comunicación de
la aplicación \textit{Tinder UD}, elaborados a partir de los casos de uso
CU-C01 a CU-C07 descritos en la Parte I.

Cada caso de uso cuenta con dos diagramas que representan la misma interacción
desde perspectivas distintas: el diagrama de secuencia destaca el orden temporal
de los mensajes, mientras que el diagrama de comunicación destaca la
organización de los objetos que participan y los mensajes que intercambian.

Los diagramas se construyeron con un enfoque de análisis, usando los tres tipos
de objeto habituales: interfaz (\textit{boundary}), controlador
(\textit{control}) y entidad de almacenamiento (base de datos).

% ---------------------------------------------------------
% 11. PARTICIPANTES Y NOTACIÓN
% ---------------------------------------------------------

\section{Participantes y Notación}

\subsection{Participantes}

\begin{center}
\begin{tabularx}{\textwidth}{|l|l|X|}
\hline
\rowcolor{gray!25}
\textbf{Participante} & \textbf{Tipo} & \textbf{Responsabilidad} \\
\hline
Estudiante & Actor & Persona que interactúa con la aplicación para alcanzar los objetivos descritos en los casos de uso. \\
\hline
Interfaz (UI) & \textit{Boundary} & Recibe las acciones del estudiante y le presenta la información y los resultados. \\
\hline
Controlador & \textit{Control} & Coordina la lógica de cada caso de uso, valida la información y decide el flujo según las condiciones. \\
\hline
Base de Datos & \textit{Database} & Almacena y consulta perfiles, intereses, conexiones, mensajes, bloqueos, reportes y cuentas. \\
\hline
\end{tabularx}
\end{center}

\subsection{Notación utilizada}

\begin{itemize}
    \item \textbf{Numeración:} los mensajes de los diagramas de secuencia se
    numeran de forma consecutiva según su orden de aparición; en los diagramas
    de comunicación se usa numeración jerárquica (1, 2, 2.1, 2.2, \dots) para
    indicar los mensajes anidados.
    \item \textbf{\texttt{alt}:} alternativas excluyentes, correspondientes a los
    escenarios alternativo y excepcional de cada caso de uso.
    \item \textbf{\texttt{opt}:} comportamiento opcional.
    \item \textbf{\texttt{loop}:} comportamiento repetitivo, por ejemplo mientras
    el estudiante continúa explorando o conversando.
    \item \textbf{\texttt{par}:} flujos que ocurren de forma concurrente.
    \item \textbf{Condiciones entre corchetes:} en los diagramas de comunicación,
    el texto entre \texttt{[ ]} indica la condición bajo la cual se envía el
    mensaje; el asterisco (\texttt{*}) indica iteración.
\end{itemize}

\newpage

% ---------------------------------------------------------
% 12. DIAGRAMAS DE SECUENCIA
% ---------------------------------------------------------

\section{Diagramas de Secuencia}

A continuación se presenta el diagrama de secuencia de cada caso de uso.

% ---------- CU-C01 ----------
\subsection{CU-C01 -- Facilitar la conexión entre estudiantes}

\Ficha{CU-C01}{Facilitar la conexión entre estudiantes}
{Estudiante, Interfaz (UI), Controlador, Base de Datos.}
{Principal: el estudiante consulta perfiles, expresa interés y, si hay interés mutuo, se establece la conexión. Alternativo: continúa explorando o modifica sus preferencias. Excepcional: no existen perfiles que coincidan con sus preferencias.}
{\texttt{loop}, \texttt{alt} anidados.}

\Diagrama{figuras/secuencia/CU-C01_secuencia.png}{fig:sec_c01}{Diagrama de secuencia de CU-C01: Facilitar la conexión entre estudiantes}

% ---------- CU-C02 ----------
\subsection{CU-C02 -- Presentarse ante otros estudiantes}

\Ficha{CU-C02}{Presentarse ante otros estudiantes}
{Estudiante, Interfaz (UI), Controlador, Base de Datos.}
{Principal: el estudiante consulta su perfil, actualiza fotografías, biografía, etiquetas, facultad y sede, y guarda la información. Excepcional: información incompleta o no válida.}
{\texttt{alt}.}

\Diagrama{figuras/secuencia/CU-C02_secuencia.png}{fig:sec_c02}{Diagrama de secuencia de CU-C02: Presentarse ante otros estudiantes}

% ---------- CU-C03 ----------
\subsection{CU-C03 -- Encontrar personas afines}

\Ficha{CU-C03}{Encontrar personas afines}
{Estudiante, Interfaz (UI), Controlador, Base de Datos, Perfiles (multiobjeto).}
{Principal: el estudiante explora perfiles y aplica criterios de intereses, facultad o sede. Alternativo: explora sin filtros específicos. Excepcional: no hay perfiles que cumplan los criterios.}
{\texttt{loop}, \texttt{alt} anidados, multiobjeto.}

\Diagrama{figuras/secuencia/CU-C03_secuencia.png}{fig:sec_c03}{Diagrama de secuencia de CU-C03: Encontrar personas afines}

% ---------- CU-C04 ----------
\subsection{CU-C04 -- Establecer una conexión}

\Ficha{CU-C04}{Establecer una conexión}
{Estudiante A, Estudiante B, Interfaz (UI), Controlador, Base de Datos.}
{Principal: ambos estudiantes expresan interés, se identifica el interés mutuo y se establece la conexión y el chat. Alternativo: no hay interés mutuo y el estudiante continúa explorando. Excepcional: el perfil deja de estar disponible.}
{\texttt{alt} anidados, \texttt{par} y nota sobre condición de carrera.}

\Diagrama{figuras/secuencia/CU-C04_secuencia.png}{fig:sec_c04}{Diagrama de secuencia de CU-C04: Establecer una conexión}

% ---------- CU-C05 ----------
\subsection{CU-C05 -- Comunicarse con una conexión}

\Ficha{CU-C05}{Comunicarse con una conexión}
{Estudiante, Conexión (otro estudiante), Interfaz (UI), Controlador, Base de Datos.}
{Principal: el estudiante selecciona una conexión, accede a la conversación y envía y recibe mensajes. Alternativo: consulta conversaciones anteriores. Excepcional: conexión no disponible, usuario bloqueado o mensaje que no puede enviarse.}
{\texttt{alt}, \texttt{loop}, \texttt{par}.}

\Diagrama{figuras/secuencia/CU-C05_secuencia.png}{fig:sec_c05}{Diagrama de secuencia de CU-C05: Comunicarse con una conexión}

% ---------- CU-C06 ----------
\subsection{CU-C06 -- Proteger la interacción}

\Ficha{CU-C06}{Proteger la interacción}
{Estudiante, Interfaz (UI), Controlador, Base de Datos.}
{Principal: el estudiante bloquea o reporta a otro usuario y confirma la acción. Alternativo: bloquea sin reportar, reporta y luego bloquea, o cancela antes de confirmar. Excepcional: usuario no disponible o reporte sin la información requerida.}
{\texttt{alt} anidados, \texttt{opt}.}

\Diagrama{figuras/secuencia/CU-C06_secuencia.png}{fig:sec_c06}{Diagrama de secuencia de CU-C06: Proteger la interacción}

% ---------- CU-C07 ----------
\subsection{CU-C07 -- Gestionar el acceso a la aplicación}

\Ficha{CU-C07}{Gestionar el acceso a la aplicación}
{Estudiante, Interfaz (UI), Controlador, Base de Datos.}
{Principal: el estudiante se registra, inicia sesión y utiliza la aplicación. Alternativo: cierra sesión, modifica su cuenta o solicita eliminarla. Excepcional: datos no válidos, cuenta ya existente o acceso no permitido.}
{\texttt{alt} anidados, \texttt{opt}.}

\Diagrama{figuras/secuencia/CU-C07_secuencia.png}{fig:sec_c07}{Diagrama de secuencia de CU-C07: Gestionar el acceso a la aplicación}

\newpage

% ---------------------------------------------------------
% 13. DIAGRAMAS DE COMUNICACIÓN
% ---------------------------------------------------------

\section{Diagramas de Comunicación}

A continuación se presenta el diagrama de comunicación de cada caso de uso. Los
mensajes conservan los mismos nombres que en los diagramas de secuencia, pero se
numeran jerárquicamente y se acompañan de las condiciones en las que se envían.

% ---------- CU-C01 ----------
\subsection{CU-C01 -- Facilitar la conexión entre estudiantes}

\Diagrama{figuras/comunicacion/CU-C01_comunicacion.png}{fig:com_c01}{Diagrama de comunicación de CU-C01: Facilitar la conexión entre estudiantes}

% ---------- CU-C02 ----------
\subsection{CU-C02 -- Presentarse ante otros estudiantes}

\Diagrama{figuras/comunicacion/CU-C02_comunicacion.png}{fig:com_c02}{Diagrama de comunicación de CU-C02: Presentarse ante otros estudiantes}

% ---------- CU-C03 ----------
\subsection{CU-C03 -- Encontrar personas afines}

\Diagrama{figuras/comunicacion/CU-C03_comunicacion.png}{fig:com_c03}{Diagrama de comunicación de CU-C03: Encontrar personas afines}

% ---------- CU-C04 ----------
\subsection{CU-C04 -- Establecer una conexión}

\Diagrama{figuras/comunicacion/CU-C04_comunicacion.png}{fig:com_c04}{Diagrama de comunicación de CU-C04: Establecer una conexión}

% ---------- CU-C05 ----------
\subsection{CU-C05 -- Comunicarse con una conexión}

\Diagrama{figuras/comunicacion/CU-C05_comunicacion.png}{fig:com_c05}{Diagrama de comunicación de CU-C05: Comunicarse con una conexión}

% ---------- CU-C06 ----------
\subsection{CU-C06 -- Proteger la interacción}

\Diagrama{figuras/comunicacion/CU-C06_comunicacion.png}{fig:com_c06}{Diagrama de comunicación de CU-C06: Proteger la interacción}

% ---------- CU-C07 ----------
\subsection{CU-C07 -- Gestionar el acceso a la aplicación}

\Diagrama{figuras/comunicacion/CU-C07_comunicacion.png}{fig:com_c07}{Diagrama de comunicación de CU-C07: Gestionar el acceso a la aplicación}

\newpage
% ---------------------------------------------------------
% 14. DIAGRAMAS DE CLASES
% ---------------------------------------------------------

\section{Diagrama de Clases}

El diagrama de clases presenta la estructura estática del sistema \textit{Tinder UD}:
las clases que representan los conceptos del dominio, sus atributos, sus operaciones
principales y las relaciones entre ellas. Se elaboró a partir de los casos de uso
CU-C01 a CU-C07 y es coherente con los diagramas de secuencia y de comunicación
presentados anteriormente.

Las clases se organizan en cinco grupos. \textbf{Usuarios:} la clase abstracta
\texttt{Usuario} agrupa los datos de acceso comunes y de ella heredan
\texttt{Estudiante}, actor principal del sistema, y \texttt{Administrador}, encargado
de revisar los reportes. \textbf{Operaciones comunes:} la interfaz genérica
\texttt{CRUD<T>} define las operaciones básicas de creación, lectura, actualización y
eliminación (\texttt{crear}, \texttt{leer}, \texttt{actualizar} y \texttt{eliminar}).
Está separada de las clases concretas para que cualquier entidad que requiera estas
operaciones pueda implementarla sin duplicar código. Actualmente la implementan
\texttt{Perfil} y \texttt{Foto}, y \texttt{Estudiante} depende de ella para gestionar
estos elementos. \textbf{Perfil:} \texttt{Perfil}, \texttt{Foto}, \texttt{Interes} y
\texttt{FiltroBusqueda} permiten que el estudiante se presente y encuentre personas
afines. \textbf{Interacción:} \texttt{Reaccion}, \texttt{Conexion} y \texttt{Mensaje}
modelan el interés entre estudiantes, la conexión que surge del interés mutuo y el
chat habilitado entre ellos. \textbf{Seguridad:} \texttt{Reporte} y \texttt{Bloqueo}
permiten controlar las interacciones no deseadas.
\Diagrama{figuras/clases/diagrama_clases.png}{fig:clases}{Diagrama de clases de Tinder UD}

\newpage

% ---------------------------------------------------------
% DIAGRAMAS DE ESTADOS
% ---------------------------------------------------------

\section{Diagramas de Estados}

Los diagramas de estados permiten representar los diferentes estados por los que
puede pasar un elemento del sistema y las condiciones que provocan las
transiciones entre ellos. Para la aplicación \textit{Tinder UD}, se elaboraron
dos diagramas de estados considerados importantes para su funcionamiento: el
estado de la cuenta del estudiante y el estado de los reportes realizados sobre
otros usuarios.

\subsection{Estado de la Cuenta}

El diagrama de estados de la cuenta representa las condiciones en las que puede
encontrarse la cuenta de un estudiante durante su uso de la aplicación. Se
definieron cuatro estados: activa, inactiva, suspendida y eliminada.

Una cuenta se encuentra activa cuando el estudiante puede acceder y utilizar las
funcionalidades de la aplicación. Si transcurren tres días sin actividad, la
cuenta pasa al estado inactivo. El estudiante puede volver a activarla iniciando
sesión correctamente. Por otra parte, un administrador puede suspender una
cuenta activa o inactiva y, posteriormente, reactivarla si lo considera
pertinente. Finalmente, una cuenta activa puede ser eliminada por solicitud del
estudiante, mientras que una cuenta suspendida puede ser eliminada por un
administrador. La eliminación corresponde al estado final de la cuenta.

\Diagrama{figuras/estados/EstadoCuenta.png}{fig:estado-cuenta}{Diagrama de estados de la cuenta del estudiante}

\subsection{Estado del Reporte}

El diagrama de estados del reporte representa el proceso que sigue un reporte
realizado por un estudiante sobre otro usuario. Este proceso contempla la
revisión del reporte, la validación de su motivo y la resolución correspondiente.

Cuando se realiza un reporte, este pasa inicialmente por una etapa de revisión
y posteriormente al estado pendiente. A continuación, se valida el motivo
indicado. Si el reporte carece de justificación o no presenta un motivo válido,
se descarta. En caso contrario, se identifica el motivo del reporte, que puede
corresponder a acoso, spam, perfil falso, contenido inapropiado u otro motivo.
Después de evaluar la situación, el reporte se marca como resuelto, finalizando
su proceso.

\Diagrama{figuras/estados/EstadoReporte.png}{fig:estado-reporte}{Diagrama de estados del reporte de un usuario}


\newpage

\newpage
% ---------------------------------------------------------
% DIAGRAMA DE ACTIVIDADES CON PARTICIONES
% ---------------------------------------------------------

\section{Diagrama de Actividades con Particiones}

El diagrama de actividades representa el flujo de trabajo de una actividad
principal del sistema \textit{Tinder UD}. Se eligió la \textbf{denuncia de un
usuario}, porque es el proceso en el que intervienen más actores y que
complementa el caso de uso CU-C06 (\textit{Proteger la interacción}) y el
diagrama de estados del reporte.

Las particiones (carriles) indican quién realiza cada actividad:
\textbf{Denunciante}, \textbf{Administrador} y \textbf{Denunciado}.

\Ficha{AD-01}{Denuncia de un usuario}
{Denunciante, Administrador y Denunciado.}
{El denunciante realiza una denuncia y el administrador la recibe y valida. Si la denuncia es válida, se bloquea la cuenta denunciada; si no lo es, se penaliza al denunciante durante 3 días. En ambos casos el administrador notifica al denunciante. El denunciado recibe una notificación y, según la decisión, el proceso termina o su cuenta se bloquea.}
{Nodo inicial, nodos de actividad, particiones, nodos de decisión, flujos de control y nodos finales.}

\subsection{Descripción del flujo}

\begin{enumerate}
    \item El \textbf{denunciante} realiza una denuncia contra otro usuario.
    \item El \textbf{administrador} recibe la denuncia y la valida.
    \item Si la validación es afirmativa (\textit{Sí}), el administrador bloquea
    la cuenta denunciada.
    \item Si la validación es negativa (\textit{No}), el administrador penaliza
    durante 3 días al denunciante.
    \item En ambos casos, el administrador notifica el resultado y el
    denunciante recibe la notificación de la denuncia, con lo que finaliza su
    flujo.
    \item El \textbf{denunciado} recibe una notificación. Si la condición
    evaluada es afirmativa (\textit{Sí}), el flujo termina; si es negativa
    (\textit{No}), su cuenta se bloquea.
\end{enumerate}

\Diagrama{figuras/actividades/diagrama_actividades.png}{fig:actividades}{Diagrama de actividades con particiones: denuncia de un usuario}

\newpage
% ---------------------------------------------------------
% DIAGRAMA DE SISTEMAS
% ---------------------------------------------------------

\section{Diagrama de Sistemas}

El diagrama de sistemas presenta la organización general de \textit{Tinder UD} en
paquetes. Muestra el sistema completo dividido en capas y la dependencia entre
ellas, e incluye el valor de inestabilidad (\textit{I}) de cada paquete, donde
$I = 0$ indica un paquete totalmente estable y $I = 1$ uno totalmente inestable.

El sistema se organiza en cuatro niveles:

\begin{itemize}
    \item \textbf{Presentación} (\texttt{UI}, $I=1.00$): contiene la interfaz con
    la que interactúa el estudiante. Depende de los controladores y nadie
    depende de ella.
    \item \textbf{Controladores} ($I=0.83$): coordinan la lógica de cada caso de
    uso y dependen de los paquetes del dominio y de la persistencia.
    \item \textbf{Dominio} ($I=0.00$): agrupa los conceptos del negocio en cinco
    paquetes: \textit{Usuarios} ($I=0.33$), \textit{Perfil} ($I=0.25$),
    \textit{Interacción} ($I=0.33$), \textit{Seguridad} ($I=0.33$) y
    \textit{Común} ($I=0.00$), que contiene la interfaz genérica
    \texttt{CRUD<T>}.
    \item \textbf{Persistencia} (\texttt{Base de Datos}, $I=0.80$): almacena la
    información del dominio.
\end{itemize}

\begin{center}
\begin{tabularx}{\textwidth}{|l|X|c|}
\hline
\rowcolor{gray!25}
\textbf{Paquete} & \textbf{Contenido} & \textbf{Inestabilidad (I)} \\
\hline
Presentación & UI & 1.00 \\
\hline
Controladores & Controladores & 0.83 \\
\hline
Usuarios & Usuario, Estudiante, Administrador & 0.33 \\
\hline
Perfil & Perfil, Foto, Interés, FiltroBúsqueda & 0.25 \\
\hline
Interacción & Reacción, Conexión, Mensaje & 0.33 \\
\hline
Seguridad & Administrar, Reporte, Bloqueo & 0.33 \\
\hline
Común & CRUD<T> & 0.00 \\
\hline
Persistencia & Base de Datos & 0.80 \\
\hline
\end{tabularx}
\end{center}

\Diagrama{figuras/sistemas/diagrama_sistemas.png}{fig:sistemas}{Diagrama de sistemas (paquetes) de Tinder UD}

\newpage
% ---------------------------------------------------------
% DIAGRAMA DE COMPONENTES
% ---------------------------------------------------------

\section{Diagrama de Componentes}

El diagrama de componentes presenta la estructura física del sistema completo en
cuatro capas y muestra qué componente soporta cada caso de uso. Los componentes
de interfaz se relacionan con los controladores mediante dependencias
(\textit{use}), y los controladores acceden a los datos mediante repositorios
conectados a la base de datos.

\begin{itemize}
    \item \textbf{Interfaz (\textit{Boundary}):} UI de Acceso a la Aplicación
    (CU-C07), UI de Administración de Reportes, UI de Protección (CU-C06),
    UI de Chat (CU-C05), UI de Conexiones (CU-C01 y CU-C04), UI de Perfil
    (CU-C02) y UI de Exploración y Filtros (CU-C03).
    \item \textbf{Control:} un controlador por cada interfaz: Acceso y Cuenta,
    Administración, Protección, Chat, Conexiones, Perfil y Exploración.
    \item \textbf{Servicios de apoyo:} Servicio de Validación de Comunidad
    Universitaria, Servicio de Notificaciones y Servicio de Almacenamiento de
    Fotos.
    \item \textbf{Acceso a datos:} Repositorio de Cuentas, Repositorio de
    Bloqueos y Reportes, Repositorio de Mensajes, Repositorio de Reacciones y
    Conexiones, y Repositorio de Perfiles e Intereses.
    \item \textbf{Base de Datos:} almacena cuentas, perfiles, fotos, intereses,
    reacciones, conexiones, mensajes, bloqueos y reportes.
\end{itemize}

Los actores que interactúan con el sistema son el \textbf{Estudiante}, que usa
la aplicación, y el \textbf{Administrador}, que revisa los reportes.

\Diagrama{figuras/componentes/diagrama_componentes.png}{fig:componentes}{Diagrama de componentes de Tinder UD}

\newpage
% ---------------------------------------------------------
% DIAGRAMA DE NODOS
% ---------------------------------------------------------

\section{Diagrama de Nodos}

El diagrama de nodos (despliegue) representa la infraestructura física sobre la
que funcionaría \textit{Tinder UD}: el hardware, el software que se ejecuta en
cada nodo y las conexiones entre ellos.

\begin{center}
\begin{tabularx}{\textwidth}{|l|X|}
\hline
\rowcolor{gray!25}
\textbf{Nodo} & \textbf{Descripción} \\
\hline
PC / Servidor principal & Equipo que aloja la aplicación. Contiene la interfaz web, el \textit{backend} (controladores, API y operaciones CRUD) y la base de datos PostgreSQL. \\
\hline
CPU & Procesador Intel Celeron G5905 de 10.ª generación, encargado del procesamiento. \\
\hline
Memoria RAM & 4 GB de memoria para el servidor. \\
\hline
Dispositivos de entrada y salida & Monitor (salida), ratón y teclado (entrada), y la interfaz de visualización y acciones. \\
\hline
Internet & Conexión de red del servidor hacia los usuarios. \\
\hline
Módem / Router & Dispositivo que conecta el servidor con Internet mediante red IP. \\
\hline
\end{tabularx}
\end{center}

Dentro del servidor, la interfaz web envía solicitudes al \textit{backend}, y
este realiza consultas SQL a la base de datos PostgreSQL. El servidor se conecta
a Internet a través del módem o router, lo que permite que los estudiantes
accedan a la aplicación.

\Diagrama{figuras/nodos/diagrama_nodos.png}{fig:nodos}{Diagrama de nodos de Tinder UD}
% =========================================================
% CONSIDERACIONES FINALES
% =========================================================

\section{Consideraciones}

Los casos de uso identificados representan los objetivos principales del estudiante
al utilizar la aplicación y no corresponden a decisiones específicas de
implementación.

Las funcionalidades como filtros, fotografías, etiquetas de intereses, ``like'',
``dislike'', generación de \textit{match} y habilitación del chat se consideran
comportamientos necesarios para alcanzar los objetivos descritos en los casos de
uso.

El caso de uso principal de la aplicación es \textit{Facilitar la conexión entre
estudiantes}, mientras que los demás casos complementan este objetivo permitiendo
que los estudiantes se presenten, encuentren personas afines, establezcan
conexiones, se comuniquen y mantengan un entorno de interacción controlado.

Los diagramas de secuencia y de comunicación describen las interacciones
necesarias para cumplir los objetivos de cada caso de uso, sin definir decisiones
de implementación como tecnologías, protocolos o estructura interna de la base de
datos. En ambos tipos de diagrama se conserva la misma separación de
responsabilidades entre interfaz, controlador y base de datos, y los escenarios
principal, alternativo y excepcional de cada caso de uso quedan representados
mediante fragmentos combinados o mensajes condicionados.

Los diagramas de estados permiten representar los cambios que pueden
experimentar dos elementos importantes de la aplicación \textit{Tinder UD}: las
cuentas de los estudiantes y los reportes realizados sobre otros usuarios. Estos
diagramas complementan los casos de uso y los diagramas de interacción, al
mostrar cómo evoluciona cada elemento según los eventos y las condiciones
definidas.

\end{document}